\chapter{The Four Ways of Reading}\index{modes!Fourfold Array|textbf}
\label{ch:four-ways}

\begin{versebox}
Etymology and intent ---\\
occasion, context, thread ---\\
four ways to read a single verse,\\
four ways to hear what's said.
\end{versebox}

\noindent Every sutta has layers. You can read the surface and miss the depth, or you can learn to ask four questions that crack any text wide open.\footnote{Pali: \textit{``Tattha katamo catubyūho hāro? `Neruttamadhippāyo'ti ayaṃ. Byañjanena suttassa neruttañca adhippāyo ca nidānañca pubbāparasandhi ca gavesitabbo.''} The mnemonic fragment is \textit{neruttam adhippāyo} (``etymology and intent''). The four elements to search for: (1) \textit{nirutti} --- the linguistic analysis; (2) \textit{adhippāya} --- the speaker's intent; (3) \textit{nidāna} --- the occasion, what prompted the statement; (4) \textit{pubbāparasandhi} --- the before-and-after connection with other teachings. The word \textit{catubyūha} literally means ``four-fold deployment,'' as in the formation of an army in four divisions.}

What are the words actually doing? What did the speaker \textit{mean}? What prompted this particular teaching at this particular moment? And how does it connect to what came before and what comes after?

These are the Four Ways. Master them, and no sutta is opaque.

\section*{The First Way: Language\index{etymology}}

The first way is linguistic analysis --- paying attention to how the words work.\footnote{Pali: \textit{``Tattha katamaṃ neruttaṃ, yā niruttipadasaṃhitā, yaṃ dhammānaṃ nāmaso ñāṇaṃ.''} \textit{Nirutti} is often translated ``etymology'' but is broader than that --- it covers the entire relationship between words and meanings, including grammar, terminology, and regional dialect.}

When a monk knows the name for a meaning and the name for a thing, the Netti says, then he can place them properly --- correctly identifying what is being talked about and how.\footnote{Pali: \textit{``Yadā hi bhikkhu atthassa ca nāmaṃ jānāti, dhammassa ca nāmaṃ jānāti, tathā tathā naṃ abhiniropeti.''} ``When a monk knows the name for the meaning and the name for the thing, he can set it up accordingly.'' Knowing what words \textit{refer to} is the foundation of all interpretation.} Such a person is called skilled in meaning, skilled in the teaching, skilled in phrasing, skilled in language, skilled in what comes before and after, skilled in presenting, skilled in the terms for past, future, present, feminine, masculine, neuter, singular, and plural.\footnote{Pali: \textit{``Ayañca vuccati atthakusalo dhammakusalo byañjanakusalo niruttikusalo pubbāparakusalo desanākusalo atītādhivacanakusalo anāgatādhivacanakusalo paccuppannādhivacanakusalo itthādhivacanakusalo purisādhivacanakusalo napuṃsakādhivacanakusalo ekādhivacanakusalo anekādhivacanakusalo.''} Fourteen linguistic competencies. This is essentially a description of a complete Pali grammarian and hermeneuticist. The list includes tense (\textit{atīta/anāgata/paccuppanna}), gender (\textit{itthī/purisa/napuṃsaka}), and number (\textit{eka/aneka}) --- the basic categories of Pali grammar, deployed here as interpretive tools.}

That sounds like a laundry list, but the point is practical: you need to know how language works before you can read what the Buddha meant. Is a word being used literally or metaphorically? Is the tense past, present, or future --- and does that matter for the meaning? Is the speaker using a regional expression that changes the sense? All of this falls under the First Way.

\section*{The Second Way: What Did the Speaker Mean\index{speaker's intent}?}

This is where reading gets interesting. Every verse has an \textit{intention} behind it --- something the Buddha was trying to accomplish. The Netti demonstrates with five examples, each time asking the same question: what was the Buddha's real point?\footnote{Pali: \textit{``Tattha katamo adhippāyo?''} ``What is the intent?'' The Netti asks this question five times, once for each example verse, and each time the answer reveals something that the surface meaning alone would miss. \textbf{A note on our approach}: the original Pali gives each verse followed by a brief statement of the Buddha's intent (\textit{ko adhippāyo?}). We have woven these into a running narrative to show the pattern more clearly, but the intent-statements are translated faithfully.}

Take this verse:

\begin{versebox}
The teaching protects the one who lives it ---\\
a great umbrella in the rainy season.\\
This is the benefit of the teaching well-practiced:\\
one who lives the teaching does not go to ruin.
\end{versebox}

What was the Buddha's intent? Not just to praise good behavior. His real point was this: those who want to escape the lower realms will become people who live the teaching.\footnote{Pali: \textit{``Dhammo have rakkhati dhammacāriṃ, chattaṃ mahantaṃ yatha vassakāle; / Esānisaṃso dhamme suciṇṇe, na duggatiṃ gacchati dhammacārī''ti.} (Jā I.10.103) \textit{``Ye apāyehi parimuccitukāmā bhavissanti, te dhammacārino bhavissantī''ti ayaṃ ettha bhagavato adhippāyo.''} The surface meaning is ``the teaching protects the practitioner.'' The deeper intent is motivational: to inspire those who fear the lower realms.} The verse is not a description --- it is an invitation.

Or this one:

\begin{versebox}
Like a thief caught at the break-in,\\
destroyed and bound by his own deed ---\\
just so, after death, in the next world,\\
destroyed and bound by their own deed.
\end{versebox}

The surface says: bad deeds have bad results. But the intent goes further. The Buddha meant: those who have deliberately done harmful deeds and accumulated them will experience unpleasant, unwanted results.\footnote{Pali: \textit{``Coro yathā sandhimukhe gahīto, sakammunā haññati bajjhate ca; / Evaṃ ayaṃ pecca pajā parattha, sakammunā haññati bajjhate cā''ti.} \textit{``Sañcetanikānaṃ katānaṃ kammānaṃ upacitānaṃ dukkhavedanīyānaṃ aniṭṭhaṃ asātaṃ vipākaṃ paccanubhavissatī''ti.''} The key word in the intent is \textit{sañcetanikānaṃ} --- ``deliberate.'' The intent narrows the verse's scope: it is not about accidental harm but about intentional, accumulated unwholesome action.} The key word is \textit{deliberate} --- the intent narrows the scope.

And this:

\begin{versebox}
Creatures that wish for happiness ---\\
whoever harms them with a rod,\\
seeking happiness for himself,\\
will find no happiness after death.
\end{versebox}

The intent? Those who want happiness will not do harmful deeds.\footnote{Pali: \textit{``Sukhakāmāni bhūtāni, yo daṇḍena vihiṃsati; / Attano sukhamesāno, pecca so na labhate sukhan''ti.} (Dhp 131--132) \textit{``Ye sukhena atthikā bhavissanti, te pāpakammaṃ na karissantī''ti.''} Again, the intent is motivational: the verse is not merely a cosmic law but an appeal to self-interest.} Again --- not a cosmic law, but an appeal to self-interest. The Buddha is saying: \textit{if you want what you say you want, then stop doing what guarantees you won't get it.}

One more:

\begin{versebox}
When one is sluggish and a glutton,\\
a sleeper, rolling on the bed ---\\
like a great hog fed on slops,\\
that fool enters the womb again and again.
\end{versebox}

A harsh image. The intent? Those who want to be free of aging and death will eat in moderation, guard their senses, devote themselves to wakefulness through the night, practice insight, cultivate wholesome states, and show respect to their companions --- elders, juniors, and those in between.\footnote{Pali: \textit{``Middhī yadā hoti mahagghaso ca, niddāyitā samparivattasāyī; / Mahāvarāhova nivāpapuṭṭho, punappunaṃ gabbhamupeti mando''ti.} (Dhp 325) The intent (\textit{adhippāya}) lists six specific practices: moderation in eating (\textit{bhojane mattaññutā}), sense-restraint (\textit{indriyesu guttadvāratā}), wakefulness (\textit{jāgariyānuyoga}), insight (\textit{vipassanā}), cultivation of wholesome states (\textit{kusalesu dhammesu}), and respect for fellow practitioners. One vivid verse conceals an entire practice regimen.} One vivid, memorable verse contains an entire practice regimen. The image of the pig is the hook; the intent is the program.

And the most famous verse of all:

\begin{versebox}
Heedfulness is the path to the deathless.\\
Heedlessness is the path to death.\\
The heedful do not die.\\
The heedless are as if already dead.
\end{versebox}

The intent: those who want to seek the deathless will dwell heedfully.\footnote{Pali: \textit{``Appamādo amatapadaṃ, pamādo maccuno padaṃ; / Appamattā na mīyanti, ye pamattā yathā matā''ti.} (Dhp 21) \textit{``Ye amatapariyesanaṃ pariyesitukāmā bhavissanti, te appamattā viharissantī''ti.''} Five verses, five intents, one pattern: the Buddha's verses are never just observations. They are always calls to action. The Second Way teaches you to hear the call.}

Five verses, five intents. And the pattern is always the same: the Buddha's verses are never just observations. They are always calls to action. The Second Way teaches you to hear the call behind the words.

\section*{The Third Way: What Prompted This\index{occasion}?}

Every teaching was given on a specific occasion, in response to a specific situation. Knowing that context can completely change what you understand.\footnote{Pali: \textit{``Tattha katamaṃ nidānaṃ?''} ``What is the occasion?'' The Pali word \textit{nidāna} means ``cause, source, occasion'' --- the circumstance that gave rise to a particular teaching.}

The Netti demonstrates with a series of exchanges. When the herdsman Dhaniya said to the Buddha:

\begin{versebox}
``A man with sons delights in sons;\\
a man with cattle delights in cattle.\\
Possessions are a person's delight ---\\
one without possessions has no delight.''
\end{versebox}

The Buddha replied:

\begin{versebox}
``A man with sons sorrows over sons;\\
a man with cattle sorrows over cattle.\\
Possessions are a person's sorrow ---\\
one without possessions has no sorrow.''
\end{versebox}

Same structure, opposite message. And the context tells you that when the Buddha says ``possessions'' here, he means external possessions --- family, cattle, property.\footnote{Pali: \textit{``Nandati puttehi puttimā, gomā gohi tatheva nandati; / Upadhī hi narassa nandanā, na hi so nandati yo nirūpadhī''ti.} (Sn 33; cf.\ SN 1.144) Dhaniya's verse, followed by the Buddha's reply: \textit{``Socati puttehi puttimā \ldots''} The key word is \textit{upadhi} (``possessions, acquisitions''). \textit{``Iminā vatthunā iminā nidānena evaṃ ñāyati `idha bhagavā bāhiraṃ pariggahaṃ upadhiṃ āhā'\,''ti.} ``By this occasion it is known that here the Blessed One meant external possessions by `acquisitions.'\,''} That is something you can only know from the context.

But when Māra the evil one hurled a great boulder from Vulture Peak at the Buddha, and the Buddha said:

\begin{versebox}
``Even if you shook the whole of\\
Vulture Peak, it would not matter ---\\
nothing can make a rightly freed\\
awakened one so much as flinch.\\[4pt]
If the sky should crack, the earth should quake,\\
and every living thing should tremble ---\\
even if a dart should shake their chest,\\
the awakened ones take no refuge in possessions.''
\end{versebox}

Here, ``possessions'' means the body. Same word, different context, different meaning.\footnote{Pali: \textit{``Yathā ca māro pāpimā gijjhakūṭā pabbatā puthusilaṃ pātesi \ldots `Neva sammāvimuttānaṃ, buddhānaṃ atthi iñjitaṃ \ldots upadhīsu tāṇaṃ na karonti buddhā'\,''ti. ``Iminā vatthunā iminā nidānena evaṃ ñāyati `idha bhagavā kāyaṃ upadhiṃ āhā'\,''ti.''} (Cf.\ SN 1.147.) By this context, the word \textit{upadhi} now means the body itself. The same word, two different contexts, two different referents.}

And elsewhere the Buddha described the body as ``diseased, impure, putrid, foul-smelling, body-dependent, oozing day and night, delighted in by fools'' --- and here too the context tells you he is talking about craving directed at the body, craving for one's own physical existence. While in another verse --- ``Uproot your own affection, like a hand plucking an autumn lily; cultivate the path of peace, the nirvana taught by the Fortunate One'' --- the occasion tells you the teaching is about abandoning that same craving.\footnote{Pali: \textit{``Āturaṃ asuciṃ pūtiṃ, duggandhaṃ dehanissitaṃ; / Paggharantaṃ divā rattiṃ, bālānaṃ abhinanditan''ti. ``Iminā vatthunā iminā nidānena evaṃ ñāyati `idha bhagavā ajjhattikavatthukāya taṇhāya pahānaṃ āhā'\,''ti.} And: \textit{``Ucchinda sinehamattano, kumudaṃ sāradikaṃva pāṇinā; / Santimaggameva brūhaya, nibbānaṃ sugatena desitan''ti.''} (Dhp 285) Two verses, same topic (craving for one's own body), different angles: one describes the problem, the other the solution. The context reveals the connection.}

The Third Way is essentially ancient literary criticism: read the context, and the text opens up.

\section*{The Fourth Way: The Thread\index{before-and-after connection}}

The fourth and most sophisticated way of reading is to connect what comes before with what comes after --- to find the thread that links one sutta to another, one verse to its explanation, one teaching to its consequence.\footnote{Pali: \textit{``Tattha katamo pubbāparasandhi.''} ``What is the before-and-after connection?'' The word \textit{sandhi} means ``junction, connection'' --- the joint where two teachings meet. This is the most sophisticated of the Four Ways because it requires knowledge of multiple texts and the ability to see patterns across them.}

The Netti shows how a verse about sensual craving connects to a verse about delusion, and how both connect to a verse about the liberated sage. Different words, different occasions --- but the \textit{thread} that runs through them is the same: craving operates through obsession, and obsession operates in three time-periods (past, present, and future), and what is called ``proliferation'' and ``obstruction'' and ``fetter'' are all different names for craving's manifestations.\footnote{Pali: \textit{``Kāmandhā jālasañchannā, taṇhāchadanachāditā; / Pamattabandhanā baddhā, macchāva kumināmukhe; / Jarāmaraṇamanventi, vaccho khīrapakova mātara''nti. Ayaṃ kāmataṇhā vuttā. Sā katamena pubbāparena yujjati? \ldots ``Ratto atthaṃ na jānāti, ratto dhammaṃ na passati; / Andhantamaṃ tadā hoti, yaṃ rāgo sahate nara''nti. Iti andhatāya ca sañchannatāya ca sāyeva taṇhā abhilapitā.''} (Ud 64; cf.\ similar phrasing at various points.) The Netti traces the thread: ``blindness'' in both verses refers to the same thing --- craving. The before-verse describes the condition; the after-verse names its mechanism. By connecting them, you see a single teaching told in two installments.}

And the thread reveals the Four Noble Truths everywhere. Craving bears fruit in three ways: results experienced in this life, in the next life, and in subsequent lives.\footnote{Pali: \textit{``Tattha pariyuṭṭhānasaṅkhārā diṭṭhadhammavedanīyā vā upapajjavedanīyā vā aparāpariyavedanīyā vā, evaṃ taṇhā tividhaṃ phalaṃ deti.''} Three time-horizons of karmic fruit. The Netti maps these onto the Buddha's statement about actions done through greed bearing results in this life, the next, or later lives. This is the ``before-and-after connection'' at work: connecting the general principle (craving bears threefold fruit) with its specific illustration (the Buddha's words on greed).} And craving is abandoned in three ways: obsession is abandoned by reflective analysis, formations are abandoned by direct seeing, and the thirty-six currents of craving are abandoned by development.\footnote{Pali: \textit{``Tattha pariyuṭṭhānaṃ paṭisaṅkhānabalena pahātabbaṃ, saṅkhārā dassanabalena, chattiṃsa taṇhāvicaritāni bhāvanābalena pahātabbānī''ti.} Three types of abandonment: (1) obsession (\textit{pariyuṭṭhāna}) by the power of reflection (\textit{paṭisaṅkhāna}); (2) formations (\textit{saṅkhārā}) by the power of seeing (\textit{dassana}); (3) the thirty-six currents of craving by the power of development (\textit{bhāvanā}). These map onto the three types of abandonment recognized in the Abhidhamma: abandonment by suppression, by cutting off, and by development.}

The total freedom from craving --- that is nirvana with a remainder. The breaking apart of the body --- that is nirvana without a remainder.\footnote{Pali: \textit{``Yā nittaṇhātā ayaṃ saupādisesā nibbānadhātu. Bhedā kāyassa ayaṃ anupādisesā nibbānadhātu.''} The two elements of nirvana: with remainder (\textit{sa-upādisesa}) = the living arahant who has destroyed craving but still has a body; without remainder (\textit{anupādisesa}) = the final passing, after which no further rebirth occurs.}

\section*{Four Types of Connection}

The before-and-after connection itself comes in four types: connection of meaning, connection of phrasing, connection of teaching-style, and connection of analysis.\footnote{Pali: \textit{``So cāyaṃ pubbāparo sandhi catubbidho atthasandhi byañjanasandhi desanāsandhi niddesasandhī''ti.} Four joints between teachings: meaning (\textit{attha}), phrasing (\textit{byañjana}), teaching-style (\textit{desanā}), and analysis (\textit{niddesa}).}

The connection of meaning has six facets: illustration, clarification, disclosure, analysis, elucidation, and designation.\footnote{Pali: \textit{``Atthasandhi chappadāni saṅkāsanā pakāsanā vivaraṇā vibhajanā uttānīkammatā paññattī''ti.}} The connection of phrasing has six facets: syllable, word, phrase, form, dialect, and exposition.\footnote{Pali: \textit{``Byañjanasandhi chappadāni akkharaṃ padaṃ byañjanaṃ ākāro nirutti niddeso''ti.} These are the building blocks of the Pali language itself: \textit{akkhara} (syllable/letter), \textit{pada} (word), \textit{byañjana} (phrase/expression), \textit{ākāra} (grammatical form), \textit{nirutti} (dialect/etymology), and \textit{niddesa} (detailed exposition). The before-and-after connection requires competence at every level of linguistic analysis.}

The connection of teaching-style is the most dramatic. The Netti illustrates it with a breathtaking passage about the meditator who depends on nothing --- not earth, not water, not fire, not wind, not any of the formless attainments, not this world or the next, not even what has been seen, heard, sensed, cognized, sought, thought about, or mentally considered. The meditator meditates without depending on \textit{any} of it --- and the whole world, including its gods, cannot fathom what is happening.\footnote{Pali: \textit{``Na ca pathaviṃ nissāya jhāyati jhāyī jhāyati ca. Na ca āpaṃ nissāya \ldots na ca tejaṃ nissāya \ldots na ca vāyuṃ nissāya \ldots Na ca ākāsānañcāyatanaṃ nissāya \ldots na ca nevasaññānāsaññāyatanaṃ nissāya \ldots na ca imaṃ lokaṃ nissāya \ldots na ca paralokaṃ nissāya jhāyati jhāyī jhāyati ca.''} This remarkable passage negates every possible support for meditation --- earth, water, fire, wind, the four formless bases, this world, the other world, and even the entire field of experience (``what is seen, heard, sensed, cognized''). Yet the meditator \textit{meditates}. The teaching-connection links this passage to Godhika's story (SN 4.23) and Vakkali's story, where Māra searched for the consciousness of the departed arahant and could not find it. The unfindability of the liberated mind \textit{is} the teaching.}

Just as Māra searched for the consciousness of the monk Godhika after his death and could not find it anywhere --- because Godhika had gone beyond proliferation, beyond craving, beyond any view to stand on --- so too Vakkali and others like him, with minds dependent on nothing, cannot be fathomed by the whole world with its gods.

And the connection of analysis is the most practical. It works in matched pairs: a mind dependent on things should be analyzed in terms of the unwholesome; a mind independent of things should be analyzed in terms of the wholesome. A dependent mind is analyzed through defilement; an independent mind through purification. Through the continuation of rebirth; through the cessation of rebirth. Through craving and ignorance; through calm and insight. Through shamelessness; through conscience. Through inattention; through mindfulness. Through unwise attention; through wise attention. Through laziness; through energy. Through faithlessness; through faith.\footnote{Pali: \textit{``Tattha katamā niddesasandhi? Nissitacittā akusalapakkhena niddisitabbā, anissitacittā kusalapakkhena niddisitabbā \ldots Nissitacittā assaddhiyena ca pamādena ca niddisitabbā, anissitacittā saddhāya ca appamādena ca niddisitabbā.''} The original lists thirteen matched pairs of analysis. We have given a representative selection here; the complete list includes: unwholesome/wholesome, defilement/purification, continuation of \textit{saṃsāra}/cessation of \textit{saṃsāra}, craving-and-ignorance/calm-and-insight, shamelessness/conscience, inattention/mindfulness, unwise attention/wise attention, laziness/energy, faithlessness/faith, not hearing the true teaching/hearing the true teaching, longing-and-ill-will/non-longing-and-non-ill-will, hindrances/the mind-liberation through dispassion and wisdom-liberation through dispassion, eternalism-and-annihilationism/the two elements of nirvana. Each pair is a lens for interpretation: dependent mind = left column; independent mind = right column.}

A dependent mind analyzed through hindrances and fetters; an independent mind analyzed through the liberation of the heart through dispassion and the liberation through wisdom. A dependent mind analyzed through the views of eternalism and annihilationism; an independent mind analyzed through the two elements of nirvana --- with remainder and without.

Every sutta, every verse, can be connected to every other through these four types of thread. That is why the Venerable Mahākaccāyana said: ``Etymology and intent.''\footnote{Pali: \textit{``Tenāha āyasmā mahākaccāyano `neruttamadhippāyo'\,''ti.''}}
